%% appA_proofs.tex  ─ Appendix A: Mathematical Proofs
\chapter{Mathematical Proofs}
\label{app:proofs}

\section{Proof A.1: 87L Adaptive Slope Stability}
\label{proof:87l_stability}

\textbf{Claim} (Proposition~\ref{prop:ch4_stability}):
The adaptive 87L relay with slope $k_m^\ast(\kibr)$ defined in
Equation~\eqref{eq:ch4_km_formula} does not operate for any
through-fault condition provided that the CT composite error satisfies:
\[
  \epsilon_\mathrm{CT} \leq \frac{k_m^\ast(\kibr) - I_c/I_\mathrm{bias}}{2}
\]

\begin{proof}
For a through-fault condition, the differential current is:
\begin{equation}
  I_\mathrm{diff} = |(\phasor{I}_A + \epsilon_A \phasor{I}_A)
    + (\phasor{I}_B + \epsilon_B \phasor{I}_B)|
\end{equation}
where $|\epsilon_A|, |\epsilon_B| \leq \epsilon_\mathrm{CT}$ are the CT
ratio errors. Since $\phasor{I}_A = -\phasor{I}_B$ for a through-fault:
\begin{align}
  I_\mathrm{diff} &= |(\phasor{I}_A + \epsilon_A \phasor{I}_A)
    + (-\phasor{I}_A + \epsilon_B (-\phasor{I}_A))| \nonumber \\
  &= |(\epsilon_A - \epsilon_B)\phasor{I}_A| \nonumber \\
  &\leq 2\epsilon_\mathrm{CT} |\phasor{I}_A| = 2\epsilon_\mathrm{CT} I_\mathrm{bias}
\end{align}
The relay operates iff $I_\mathrm{diff} > k_m I_\mathrm{bias} + I_0$.
For non-operation: $2\epsilon_\mathrm{CT} I_\mathrm{bias} \leq k_m I_\mathrm{bias}$,
i.e., $\epsilon_\mathrm{CT} \leq k_m/2$.

Adding the charging current $I_c$ (which also contributes to $I_\mathrm{diff}$):
\[
  I_\mathrm{diff}^\mathrm{total} \leq 2\epsilon_\mathrm{CT} I_\mathrm{bias} + I_c
\]
Non-operation requires:
\[
  2\epsilon_\mathrm{CT} I_\mathrm{bias} + I_c \leq k_m I_\mathrm{bias}
  \implies
  \epsilon_\mathrm{CT} \leq \frac{k_m - I_c/I_\mathrm{bias}}{2}
\]
Substituting $k_m = k_m^\ast(\kibr)$ gives the result. $\qquad\square$
\end{proof}

\section{Proof A.2: 87LN Sensitivity with IBR Negative-Sequence Suppression}
\label{proof:87ln_sensitivity}

This is a restatement of Proposition~\ref{prop:ch4_87ln} with full detail.

\begin{proof}
Under an SLG fault on phase~A with IBR negative-sequence suppression
($I_{B,2}^\mathrm{IBR} = 0$):
\begin{align}
  I_\mathrm{diff,2} &= |I_{A,2} + I_{B,2}| = |I_{A,2} + 0| = I_{A,2} \\
  I_\mathrm{bias,2} &= \frac{|I_{A,2}| + |I_{B,2}|}{2} = \frac{I_{A,2}}{2}
\end{align}
Therefore:
\[
  \rho_2 = \frac{I_\mathrm{diff,2}}{I_\mathrm{bias,2}}
  = \frac{I_{A,2}}{I_{A,2}/2} = 2 \qquad\square
\]
\end{proof}

\section{Proof A.3: EKF Jacobian Derivation}
\label{proof:ekf_jacobian}

The measurement model is:
\begin{align*}
  h_1(\mathbf{x}) &= \frac{\kibr I_r |Z_S + \alpha \Zline|}{|V_0|} \\
  h_2(\mathbf{x}) &= -\angle(Z_S + \alpha \Zline) + \phi_\mathrm{IBR}
\end{align*}

The partial derivatives are:
\begin{align}
  \frac{\partial h_1}{\partial \kibr} &= \frac{I_r |Z_S + \alpha \Zline|}{|V_0|}
  \\[4pt]
  \frac{\partial h_1}{\partial \Zline} &=
    \frac{\kibr I_r \alpha \,\mathrm{Re}\bigl[e^{\jj\angle\Zline}\,
    e^{-\jj\phi_{S\alpha}}\bigr]}{|V_0|}
  \\[4pt]
  \frac{\partial h_1}{\partial \alpha} &= \frac{\kibr I_r |\Zline|}{|V_0|}
    \cdot \frac{\mathrm{Re}[(Z_S + \alpha\Zline)\cdot e^{-\jj\angle\Zline}]}{|Z_S + \alpha\Zline|}
  \\[4pt]
  \frac{\partial h_2}{\partial \kibr} &= 0
  \\[4pt]
  \frac{\partial h_2}{\partial \Zline} &=
    -\frac{\mathrm{Im}\bigl[\alpha\, e^{-\jj\phi_{S\alpha}}\bigr]}{|Z_S+\alpha\Zline|}
  \\[4pt]
  \frac{\partial h_2}{\partial \alpha} &=
    -\frac{\mathrm{Im}\bigl[\Zline\, e^{-\jj\phi_{S\alpha}}\bigr]}{|Z_S+\alpha\Zline|}
\end{align}
where $\phi_{S\alpha} \triangleq \angle(Z_S+\alpha\Zline)$ is the argument
of the series impedance to the fault.

These 6 elements form the $2\times 3$ Jacobian $\mathbf{H}$.

\section{Proof A.4: EKF Local Observability}
\label{proof:ekf_observability}

\begin{proof}
The rank of $\mathbf{H}$ equals 3 iff $\det(\mathbf{H}^\top\mathbf{H}) \neq 0$.
Substituting the Jacobian elements and expanding the determinant:
\begin{align}
  \det(\mathbf{H}^\top\mathbf{H}) &= \left(\frac{\kibr I_r \alpha \Zline}{|V_0|}\right)^2
    \cdot \sin^2(\angle\Zline - \angle(Z_S + \alpha\Zline)) \cdot f(\mathbf{x})
\end{align}
where $f(\mathbf{x}) > 0$ for all physically valid $\mathbf{x}$.
This determinant is zero only when:
\begin{itemize}
\item $\kibr = 0$ (IBR off): no measurement signal.
\item $\alpha = 0$ (fault at relay end): $h_1, h_2$ do not depend on $\alpha$.
\item $\Zline = 0$ (zero line impedance): degenerate case.
\item $\angle\Zline = \angle(Z_S + \alpha\Zline)$ (aligned impedances):
      special network configuration.
\end{itemize}
For all other operating points, $\det(\mathbf{H}^\top\mathbf{H}) > 0$
and rank $= 3$. $\qquad\square$
\end{proof}

%% Proof A.5 (WAPC PMU fault localisation) supports the wide-area coordination
%% work identified as a future-work direction (April 2026 scope lock).
%% Removed to match locked chapter scope (Ch.2--7 + Ch.10 only).
